\section{Dzienniczek pacjenta / ePRO}

% SLIDE ID: sec04-goal
\BlockGoalFrame{\SlideID{sec04-goal}}{
  \item pokazać, jak dane mogą być zbierane bezpośrednio od pacjenta,
  \item zrozumieć wpływ ePRO na ciągłość i jakość danych,
  \item omówić korzyści, ograniczenia i ryzyka tego modelu pracy.
}

% SLIDE ID: sec04-what-is-epro
\begin{frame}{Czym jest ePRO}
  \SlideID{sec04-what-is-epro}
  \begin{DefinitionBox}
    ePRO to elektroniczne zbieranie danych zgłaszanych bezpośrednio przez pacjenta, bez pośredniego przepisywania ich przez personel ośrodka.
  \end{DefinitionBox}

  \vspace{0.35cm}

  \begin{itemize}
    \item wpis odbywa się w aplikacji lub dzienniczku elektronicznym,
    \item dane pojawiają się szybciej niż przy klasycznym kontakcie wizytowym,
    \item system może wspierać regularność wpisów i ich kompletność.
  \end{itemize}
\end{frame}

% SLIDE ID: sec04-diary
\begin{frame}{Czym jest dzienniczek pacjenta}
  \SlideID{sec04-diary}
  \begin{itemize}
    \item to narzędzie do codziennych lub cyklicznych wpisów pacjenta,
    \item może działać na telefonie, tablecie albo w przeglądarce,
    \item zwykle służy do zbierania objawów, przyjmowania leku i pomiarów domowych.
  \end{itemize}
\end{frame}

% SLIDE ID: sec04-history-overview
\begin{frame}{Historia rozwoju ePRO}
  \SlideID{sec04-history-overview}
  \begin{center}
    \Large
    Wczesne systemy elektroniczne $\rightarrow$ telefoniczne \texttt{IVRS} $\rightarrow$ urządzenia dedykowane $\rightarrow$ \texttt{BYOD} i badania zdecentralizowane
  \end{center}

  \vspace{0.45cm}

  \begin{itemize}
    \item Historia \texttt{ePRO} to przejście od papierowych dzienniczków, które łatwo było uzupełniać po fakcie, do danych zgłaszanych bliżej rzeczywistego momentu objawu.
    \item Najważniejsza zmiana nie dotyczyła tylko nośnika, ale wiarygodności czasu wpisu, ograniczenia recall bias i większej ciągłości obserwacji pacjenta.
    \item Z czasem \texttt{ePRO} przestało być dodatkiem do badania i zaczęło pełnić rolę ważnego źródła danych klinicznych oraz operacyjnych.
  \end{itemize}
\end{frame}

% SLIDE ID: sec04-history-early-computing
\begin{frame}{1. Wczesne systemy elektroniczne: lata 80. i 90.}
  \SlideID{sec04-history-early-computing}
  \begin{itemize}
    \item Już w latach 80. i 90. badacze sprawdzali, czy papierowe dzienniczki można zastąpić prostymi systemami elektronicznymi.
    \item Były to często rozwiązania cięższe technologicznie niż dziś: proste terminale, ekrany komputerowe lub dedykowane stanowiska zbierania danych.
    \item Początkowo obawiano się, że starsi pacjenci lub osoby mniej techniczne będą odrzucać taki model pracy, ale pilotaże często pokazywały dobrą akceptację.
    \item To był pierwszy krok do myślenia, że pacjent może raportować dane bezpośrednio do systemu, a nie tylko do papierowego formularza oddawanego na wizycie.
  \end{itemize}
\end{frame}

% SLIDE ID: sec04-history-ivrs
\begin{frame}{2. Telefoniczne ePRO i IVRS: koniec lat 90. i lata 2000}
  \SlideID{sec04-history-ivrs}
  \begin{itemize}
    \item Zanim smartfony i aplikacje stały się powszechne, ważną rolę odegrały systemy \texttt{IVRS}, w których pacjent dzwonił na numer i odpowiadał na pytania z klawiatury telefonu.
    \item Taki model ograniczał klasyczne \texttt{back-filling}, czyli uzupełnianie kilku dni dzienniczka tuż przed wizytą kontrolną.
    \item Dzięki znacznikom czasu dane były bliższe rzeczywistemu momentowi zdarzenia, co zmniejszało ryzyko zniekształceń pamięciowych.
    \item \texttt{IVRS} pokazał też, że raportowanie pacjenta może działać poza ośrodkiem i nie musi wymagać stałego kontaktu z personelem badania.
  \end{itemize}
\end{frame}

% SLIDE ID: sec04-history-handheld
\begin{frame}{3. Dedykowane urządzenia handheld: lata 2000 i 2010}
  \SlideID{sec04-history-handheld}
  \begin{itemize}
    \item W kolejnej fazie sponsorzy zaczęli przekazywać pacjentom dedykowane urządzenia lub \texttt{PDA}, używane wyłącznie do dzienniczków objawów i kwestionariuszy jakości życia.
    \item Ten model dawał większą kontrolę nad wersją oprogramowania, konfiguracją pytań, przypomnieniami i bezpieczeństwem danych.
    \item Jednocześnie oznaczał większy narzut logistyczny: urządzenia trzeba było przygotować, dystrybuować, szkolić pacjentów i serwisować.
    \item W tym okresie \texttt{ePRO} umocniło się jako poważne narzędzie zbierania danych, a nie eksperyment technologiczny.
  \end{itemize}
\end{frame}

% SLIDE ID: sec04-history-modern
\begin{frame}{4. BYOD, pandemia i model zdalny: lata 2020 do dziś}
  \SlideID{sec04-history-modern}
  \begin{itemize}
    \item Współczesny rynek mocno przesunął się w stronę \texttt{BYOD} (\texttt{Bring Your Own Device}), czyli użycia własnego telefonu pacjenta lub bezpiecznego portalu webowego.
    \item W czasie pandemii rozwiązania zdalne stały się kluczowe dla utrzymania ciągłości badań i rozwoju modeli hybrydowych oraz zdecentralizowanych.
    \item Nowoczesne platformy \texttt{ePRO} integrują przypomnienia, alerty, monitoring compliance, czasem także wearables i inne elementy \texttt{eCOA}.
    \item W praktyce \texttt{ePRO} stało się jednym z filarów patient-centric trial design: mniej wizyt tylko po to, by „przepisać” objawy, więcej bieżącej obserwacji między wizytami.
  \end{itemize}
\end{frame}

% SLIDE ID: sec04-history-clinical-value
\begin{frame}{Kamień milowy: kiedy ePRO pokazało wartość kliniczną}
  \SlideID{sec04-history-clinical-value}
  \begin{itemize}
    \item Ważnym momentem była praca zespołu \texttt{Ethan Basch} opublikowana w \texttt{Journal of Clinical Oncology} w 2016 r.
    \item W badaniu pacjenci onkologiczni raportowali objawy przez portal webowy, a zespół otrzymywał sygnały o problemach między wizytami.
    \item Pokazano nie tylko lepszą jakość życia i mniej wizyt ratunkowych, ale także wydłużenie przeżycia całkowitego w porównaniu z grupą standardowej opieki.
    \item To był ważny sygnał, że \texttt{ePRO} nie służy wyłącznie do „ładniejszego zbierania danych”, ale może realnie zmieniać przebieg opieki i bezpieczeństwo pacjenta.
  \end{itemize}
\end{frame}

% SLIDE ID: sec04-differences
\begin{frame}{ePRO, eCOA i klasyczny CRF}
  \SlideID{sec04-differences}
  \begin{columns}[T,onlytextwidth]
    \column{0.31\textwidth}
    \begin{beamercolorbox}[wd=\linewidth,sep=0.8em,rounded=true]{block body}
      \textbf{\textcolor{UAFMTeal}{ePRO}}\par\medskip
      dane raportowane bezpośrednio przez pacjenta
    \end{beamercolorbox}

    \column{0.31\textwidth}
    \begin{beamercolorbox}[wd=\linewidth,sep=0.8em,rounded=true]{block body}
      \textbf{\textcolor{UAFMDarkBlue}{eCOA}}\par\medskip
      szersza kategoria elektronicznych ocen klinicznych
    \end{beamercolorbox}

    \column{0.31\textwidth}
    \begin{beamercolorbox}[wd=\linewidth,sep=0.8em,rounded=true]{block body}
      \textbf{\textcolor{UAFMOrange}{CRF}}\par\medskip
      dane wpisywane przez personel badawczy
    \end{beamercolorbox}
  \end{columns}
\end{frame}

% SLIDE ID: sec04-patient-source
\begin{frame}{Pacjent jako źródło danych}
  \SlideID{sec04-patient-source}
  \begin{itemize}
    \item objawy i samopoczucie między wizytami,
    \item jakość życia,
    \item przyjmowanie leku,
    \item działania niepożądane,
    \item pomiary domowe, np. ciśnienie lub glikemia.
  \end{itemize}
\end{frame}

% SLIDE ID: sec04-compliance
\begin{frame}{Compliance pacjenta}
  \SlideID{sec04-compliance}
  \begin{itemize}
    \item system może mierzyć regularność wpisów,
    \item wysoki compliance zwiększa użyteczność danych,
    \item niski compliance oznacza luki, opóźnienia i większe ryzyko błędnej interpretacji.
  \end{itemize}
\end{frame}

% SLIDE ID: sec04-reminders
\begin{frame}{Przypomnienia i brakujące dane}
  \SlideID{sec04-reminders}
  \begin{itemize}
    \item przypomnienia pomagają utrzymać rytm wpisów,
    \item system może oznaczać niewypełnione wpisy lub opóźnienia,
    \item brakujące dane nadal pozostają problemem klinicznym i analitycznym.
  \end{itemize}
\end{frame}

% SLIDE ID: sec04-integration
\begin{frame}{Integracja z EDC}
  \SlideID{sec04-integration}
  \begin{itemize}
    \item dane z dzienniczka mogą trafiać do EDC automatycznie,
    \item ważne jest mapowanie pacjenta, wizyty i właściwych formularzy,
    \item dobra integracja zmniejsza ręczne przepisywanie i ryzyko błędów.
  \end{itemize}
\end{frame}

% SLIDE ID: sec04-risks
\begin{frame}{Co może pójść źle?}
  \SlideID{sec04-risks}
  \begin{itemize}
    \item słaba jakość wpisów lub niepełne odpowiedzi,
    \item przeciążenie pacjenta liczbą pytań,
    \item problemy z prywatnością i bezpieczeństwem dostępu,
    \item trudności techniczne po stronie urządzenia lub aplikacji.
  \end{itemize}
\end{frame}

% SLIDE ID: sec04-demo
\DemoFrame{Dzienniczek pacjenta}{
  interfejs pacjenta, harmonogram wpisów, przypomnienia, widok brakujących danych oraz przykład przepływu informacji do systemu centralnego.
}

% SLIDE ID: sec04-quality
\begin{frame}{Wpływ ePRO na jakość i ciągłość danych}
  \SlideID{sec04-quality}
  \begin{itemize}
    \item dane są zbierane bliżej momentu wystąpienia objawu lub zdarzenia,
    \item maleje ryzyko odtwarzania informacji z pamięci podczas wizyty,
    \item znaczniki czasu utrudniają zbiorcze uzupełnianie wpisów po fakcie,
    \item jednocześnie rośnie znaczenie użyteczności narzędzia, motywacji pacjenta i jakości projektu kwestionariusza.
  \end{itemize}
\end{frame}

% SLIDE ID: sec04-case
\CaseStudyFrame{Pacjent zgłasza objaw w dzienniczku — co dzieje się dalej?}{
  Pacjent wpisuje w dzienniczku nasilony ból głowy i zaznacza, że objaw pojawił się po przyjęciu dawki badanego leku.
}{
  System może zapisać wpis, oznaczyć go jako istotny klinicznie, przekazać dane do dalszego przeglądu i uruchomić proces oceny po stronie zespołu badania.
}

% SLIDE ID: sec04-sources
\begin{frame}{Źródła do tej sekcji}
  \SlideID{sec04-sources}
  \footnotesize
  \begin{itemize}
    \item Applied Clinical Trials, \textit{ePRO and Ensuring Data Integrity}.
    \item FDA, \textit{Patient-Reported Outcome Measures: Use in Medical Product Development to Support Labeling Claims}, December 2009.
    \item Basch E. et al., \textit{Symptom Monitoring With Patient-Reported Outcomes During Routine Cancer Treatment: A Randomized Controlled Trial}, \texttt{JCO}, 2016.
    \item Stone A.A. et al. i późniejsza literatura o elektronicznych dzienniczkach, recall bias i back-filling.
    \item Nowsze rekomendacje dotyczące projektowania elektronicznych dzienniczków zdarzeń i \texttt{ePRO} w badaniach klinicznych.
  \end{itemize}
\end{frame}
